\section{Conclusion}
\label{sec:conclusion}

In this work, we have presented a comprehensive, technically rigorous, and honest systems-ML comparative study of \textbf{Zero-Shot Calibration-Free Model Merging}. We have evaluated and compared two distinct dynamic serving paradigms under streaming conditions: (1) \textbf{Zero-Shot Expert Entropy Routing (EER)} and (2) \textbf{Entropy-Pseudo-Labeled Online Centroid Adaptation (EPL-OCA)}. Both methods eliminate the costly and restrictive operational requirement of offline labeled calibration datasets ($|\mathcal{C}_k|=64$) that bottlenecked previous dynamic merging frameworks.

Empirically, our direct-routing paradigm (\textbf{EER}) achieves outstanding results, delivering \textbf{71.38\%} Joint Mean accuracy on heterogeneous shuffled streams and \textbf{71.18\%} under extreme continuous covariate drift ($d=0.45$). This outpaces the supervised SOTA baseline (SPS-ZCA, 66.76\%) by \textbf{+4.62\%} absolute accuracy completely calibration-free, demonstrating excellent domain resilience. Conversely, we honestly report that the centroid-routing paradigm (\textbf{EPL-OCA}) only achieves \textbf{49.88\%} (and \textbf{49.78\%} under drift) due to the \textbf{Representational Sparsity Paradox}. Orthogonal class representations within task subspaces introduce high spatial sparseness and noise in centroid-based routing, leading to ensembling degradation.

From a systems perspective, we have head-on analyzed the physical serving complexity under post-Layer-4 activation divergence, detailing a realistic FLOP complexity of $0.25 + 0.75K$ passes. We propose \textbf{Amortized Pseudo-Labeling} as a highly practical systems-level mitigation to bring serving latency back to a lightweight $\approx 1.3\times$ passes, and outline an explicit engineering roadmap to resolve the representational sparsity bottleneck of centroid-based ensembling. By establishing robust, calibration-free serving boundaries, our work bridges the gap between theoretical model ensembling and resource-constrained edge deployment.

Moving forward, several exciting research directions emerge from this comparative study. First, while we have successfully implemented and evaluated the Streaming K-Means baseline, exploring more advanced, non-parametric online clustering models could offer improved resilience against non-isotropic noise. Streaming K-Means implicitly assumes isotropic, spherical task clusters of equal variance, which frequently fails under the directional, non-isotropic representation noise found in actual deep backbones. Non-parametric models like streaming Dirichlet Process Mixture Models (DPMs) resolve this by explicitly modeling cluster covariance matrices (e.g., via conjugate Normal-Inverse-Wishart priors), allowing them to dynamically adapt to ellipsoidal manifolds and arbitrary noise orientations without requiring a pre-specified $K$. Alternatively, streaming DBSCAN variants can group representations based on dense connectivity rather than distance to a single central prototype, enabling the extraction of arbitrarily shaped non-Gaussian manifolds while naturally filtering out out-of-distribution background noise as outlier coordinates. Second, exploring periodic cluster re-alignment and centroid re-indexing under extremely long-running non-linear drift would provide valuable insights into maintaining long-term serving stability. We leave these systems-ML investigations for future work.
\vfil
